%! AP Statistics — Unit 7, Lesson 7.2 — Group Activity (Answer Key)
\documentclass[10pt]{article}
\usepackage{apstats-article}
\usepackage{apstats-key}
\newcommand{\TallMath}[1]{$\displaystyle #1\rule[-1.4em]{0pt}{3.2em}$}

\begin{document}

\pageheader{Unit 7, Lesson 7.2}{Group Activity --- Answer Key}
\namepartnerperiod

\begin{tcolorbox}[colback=white, colframe=black!40, title=\textbf{Tier R --- Remediate},
    fonttitle=\bfseries, arc=2mm, left=3mm, right=3mm, top=2mm, bottom=2mm]

\textbf{Sleep Study.} $n = 25$, $\bar x = 6.8$ hr, $s = 1.2$ hr, 90\% CI.

\textbf{Step 1:} \ansline{$\mu = $ true mean nightly sleep duration (hours) for all students at the school.}

\textbf{Step 2:} Conditions pre-verified. \checkmark

\textbf{Step 3:}

$df = $ \ans{24} \quad $t^* = $ \ans{1.711}

$SE = 1.2/\sqrt{25} = $ \ans{0.240} hr

$ME = 1.711 \times 0.240 = $ \ans{0.411} hr

Interval: $6.8 \pm 0.411 = ($ \ans{6.389} $,$ \ans{7.211} $)$ hours

\textbf{Step 4:} \ansline{We are 90\% confident that the interval (6.39, 7.21) hours captures the true mean nightly sleep duration for all students at this school.}

\begin{teachernote}
    $t^*$ for $df = 24$, 90\% CI is 1.711 (from table). Some calculators give 1.7109.
\end{teachernote}

\end{tcolorbox}

\begin{tcolorbox}[colback=white, colframe=black!40, title=\textbf{Tier A --- Approaching Proficiency},
    fonttitle=\bfseries, arc=2mm, left=3mm, right=3mm, top=2mm, bottom=2mm]

\textbf{Commute Times.} $n = 36$, $\bar x = 28.4$ min, $s = 9.6$ min, 95\% CI.

\textbf{Step 1:} \ansline{$\mu = $ true mean commute time (minutes) for all workers in the downtown district.}

\textbf{Step 2:}
\ansline{Random: random sample of workers. \checkmark \quad 10\%: $36 \le 10\%$ of all downtown workers (assuming $N \ge 360$). \checkmark \quad Large Sample: $n = 36 \ge 30$, so CLT applies. \checkmark}

\textbf{Step 3:}

$df = $ \ans{35} \quad $t^* = $ \ans{2.030} \quad $SE = 9.6/\sqrt{36} = $ \ans{1.60} min

Interval: $28.4 \pm 2.030 \times 1.60 = 28.4 \pm 3.25 = ($ \ans{25.15} $,$ \ans{31.65} $)$ min

\textbf{Step 4:} \ansline{We are 95\% confident that the interval (25.15, 31.65) minutes captures the true mean commute time for all workers in the downtown district.}

\textbf{Bonus:} \ansline{Width $= 2 \times 3.25 = 6.50$ min. To halve the margin of error, quadruple the sample size: from $n = 36$ to $n = 144$.}

\end{tcolorbox}

\begin{tcolorbox}[colback=white, colframe=black!40, title=\textbf{Tier E --- Extension},
    fonttitle=\bfseries, arc=2mm, left=3mm, right=3mm, top=2mm, bottom=2mm]

\textbf{Matched-Pairs: Step-Count Intervention.}

Differences: 1.2, 0.8, $-$0.3, 2.1, 1.5, 0.6, 1.8, $-$0.1, 2.4, 0.9, 1.3, 1.7

$\bar d = $ \ans{$13.9/12 \approx 1.158$} thousands \quad $s_d \approx $ \ans{0.784} thousands \quad $n = $ \ans{12}

\textbf{Step 2:} \ansline{Random: random sample of participants (or assume randomized assignment to before/after design). \checkmark \quad 10\%: $12 \le 10\%$ of all potential participants. \checkmark \quad Large Sample: $n = 12 < 30$; must check for strong skewness or outliers in the differences. With no extreme values visible, condition is approximately met. \checkmark}

\textbf{Step 3:} $df = $ \ans{11} \quad $t^* = $ \ans{2.201}

Interval: $1.158 \pm 2.201 \times 0.784/\sqrt{12} = 1.158 \pm 2.201 \times 0.226
= 1.158 \pm 0.498 = ($ \ans{0.660} $,$ \ans{1.656} $)$ thousands of steps

\textbf{Step 4 and claim:} \ansline{We are 95\% confident that the true mean change in daily step count for participants is between 0.660 and 1.656 thousand steps. Because the entire interval lies above 0, the data provide convincing evidence that the intervention increased step counts on average.}

\begin{teachernote}
    $\bar d$: sum $= 1.2+0.8-0.3+2.1+1.5+0.6+1.8-0.1+2.4+0.9+1.3+1.7 = 13.9$;
    $\bar d = 13.9/12 \approx 1.1583$.
    $s_d$: compute via calculator; approximately 0.784.
    $t^*$ for $df = 11$, 95\%: 2.201.
\end{teachernote}

\end{tcolorbox}

\end{document}
